\documentclass[11pt]{article}
\usepackage[margin=1in]{geometry}
\usepackage{amsmath,amssymb,bm}
\usepackage{mathtools}
\usepackage{mathrsfs}
\usepackage{booktabs}
\usepackage[colorlinks=true,linkcolor=blue,urlcolor=blue]{hyperref}
\usepackage{enumitem}

\newcommand{\Rh}{\hat{\mathbf{R}}}
\newcommand{\Rhh}[1]{\hat{R}_{#1}}
\newcommand{\one}{\mathbf{1}}
\newcommand{\eps}{\bm{\epsilon}}
\newcommand{\half}{\tfrac{1}{2}}
\newcommand{\third}{\tfrac{1}{3}}

\title{\vspace{-2em}Regularized different-sized RPY pair mobilities\\
in the Kim \& Karrila irreducible basis}
\author{Translation--translation, translation--rotation, and translation--dipole couplings}
\date{}

\begin{document}
\maketitle
\vspace{-2.5em}

\section*{Purpose and conventions}

These notes transcribe the different-sized regularized RPY mobility tensors of
Zuk, Wajnryb, Mizerski \& Szymczak (\emph{J.\ Fluid Mech.}\ \textbf{741}, R5, 2014) into
the Kim \& Karrila (K\&K) scalar mobility functions, suitable for a Stokesian-dynamics
code that stores pair mobilities in the irreducible $(x,y)$ form.  For each pair $(i,j)$ we
write $\mathbf{R}_{ij}=\mathbf{R}_i-\mathbf{R}_j$, $R=|\mathbf{R}_{ij}|$,
$\Rh=\mathbf{R}_{ij}/R$; $a_i,a_j$ are the two radii and $a_>,a_<$ the larger/smaller of the
pair.  All expressions are the open-boundary (free-space) propagator forms; $\eta$ is the
solvent viscosity, and $\Theta$ is the Heaviside step function.

Each coupling has three branches: the non-overlapping far-field form
($R>a_i+a_j$), the overlapping regularized form ($a_>-a_<<R\le a_i+a_j$), and the
complete-overlap form ($R\le a_>-a_<$, one sphere inside the other).

\paragraph{Basis tensors.}
The three couplings are decomposed as
\begin{align}
\bm{\mu}^{tt}_{ij} &= x^a_{ij}\,\Rh\Rh + y^a_{ij}\,(\one-\Rh\Rh),\\[2pt]
\bm{\mu}^{rt}_{ij} &= y^b_{ij}\,\big(\eps\cdot\Rh\big),
\qquad [\eps\cdot\Rh]_{\alpha\beta}=\epsilon_{\alpha\beta\gamma}\Rhh{\gamma},\\[2pt]
\mu^{td}_{ij,\alpha\beta\gamma} &=
x^g_{ij}\Big(\Rhh{\beta}\Rhh{\gamma}-\third\delta_{\beta\gamma}\Big)\Rhh{\alpha}
+ y^g_{ij}\Big(\delta_{\alpha\beta}\Rhh{\gamma}+\delta_{\alpha\gamma}\Rhh{\beta}
-2\Rhh{\alpha}\Rhh{\beta}\Rhh{\gamma}\Big).
\label{eq:tdbasis}
\end{align}
In \eqref{eq:tdbasis} the first index $\alpha$ is the translational (velocity) index and the
pair $(\beta,\gamma)$ contracts with the symmetric traceless rate-of-strain $E_{\beta\gamma}$
(equivalently the stresslet).  Both basis tensors are symmetric and traceless in
$(\beta,\gamma)$.  This is the standard K\&K \emph{gt}/$\tilde{g}$ decomposition; see the
remark at the end regarding the normalization of the $y^g$ basis tensor.

\section{Translation--translation coupling \texorpdfstring{$\bm{\mu}^{tt}$}{mu-tt}}

Zuk et al.\ (3.1), in the $\{\one,\Rh\Rh\}$ form:
\begin{equation}
\bm{\mu}^{tt}_{ij}=
\begin{cases}
\dfrac{1}{8\pi\eta R}\!\left[\left(1+\dfrac{a_i^2+a_j^2}{3R^2}\right)\one
+\left(1-\dfrac{a_i^2+a_j^2}{R^2}\right)\Rh\Rh\right], & a_i+a_j<R,\\[14pt]
\dfrac{1}{6\pi\eta a_i a_j}\!\left[\begin{aligned}
&\tfrac{16R^3(a_i+a_j)-\big((a_i-a_j)^2+3R^2\big)^2}{32R^3}\,\one\\[2pt]
&+\tfrac{3\big((a_i-a_j)^2-R^2\big)^2}{32R^3}\,\Rh\Rh
\end{aligned}\right], & a_>-a_<<R\le a_i+a_j,\\[20pt]
\dfrac{1}{6\pi\eta a_>}\,\one, & R\le a_>-a_<.
\end{cases}
\end{equation}
Writing $\bm{\mu}^{tt}=x^a\Rh\Rh+y^a(\one-\Rh\Rh)$ (so $x^a$ is the longitudinal /
$\Rh\Rh$ eigenvalue and $y^a$ the transverse one), one has $y^a=(\text{coeff of }\one)$ and
$x^a=(\text{coeff of }\one)+(\text{coeff of }\Rh\Rh)$.  For the non-overlapping branch
($a_i+a_j<R$):
\begin{equation}
\boxed{\;
x^a_{ij}=\frac{1}{4\pi\eta R}\left(1-\frac{a_i^2+a_j^2}{3R^2}\right),
\qquad
y^a_{ij}=\frac{1}{8\pi\eta R}\left(1+\frac{a_i^2+a_j^2}{3R^2}\right).\;}
\end{equation}
For the overlapping branch ($a_>-a_<<R\le a_i+a_j$):
\begin{equation}
\boxed{\;
\begin{aligned}
x^a_{ij}&=\frac{16R^3(a_i+a_j)+2(a_i-a_j)^4-12(a_i-a_j)^2R^2-6R^4}{192\pi\eta\, a_i a_j\,R^3},\\[6pt]
y^a_{ij}&=\frac{16R^3(a_i+a_j)-\big((a_i-a_j)^2+3R^2\big)^2}{192\pi\eta\, a_i a_j\,R^3}.
\end{aligned}\;}
\end{equation}
For the complete-overlap branch ($R\le a_>-a_<$): $x^a_{ij}=y^a_{ij}=1/(6\pi\eta a_>)$.

\section{Translation--rotation coupling \texorpdfstring{$\bm{\mu}^{rt}$}{mu-rt}}

Zuk et al.\ (3.5) gives $\bm{\mu}^{rt}_{ij}=y^b_{ij}\,(\eps\cdot\Rh)$ with
\begin{equation}
\boxed{\;
y^b_{ij}=
\begin{cases}
\dfrac{1}{8\pi\eta R^2}, & a_i+a_j<R,\\[12pt]
\dfrac{(a_i-a_j+R)^2\big(a_j^2+2a_j(a_i+R)-3(a_i-R)^2\big)}{128\pi\eta\, a_i^3 a_j\,R^2},
& a_>-a_<<R\le a_i+a_j,\\[12pt]
\Theta(a_i-a_j)\,\dfrac{R}{8\pi\eta a_i^3}, & R\le a_>-a_<.
\end{cases}\;}
\end{equation}
The complete-overlap branch uses $\zeta^{rr}_i=8\pi\eta a_i^3$.  The
translation--rotation partner $\bm{\mu}^{tr}_{ij}$ has the identical form with $a_i$ and
$a_j$ interchanged (Zuk et al.\ note that $\bm{\mu}^{tr}$ is \emph{not} simply the transpose
of $\bm{\mu}^{rt}$, correcting a misprint in Wajnryb et al.\ 2013).

\section{Translation--dipole coupling \texorpdfstring{$\bm{\mu}^{td}\to(x^g,y^g)$}{mu-td to (xg,yg)}}

\subsection*{The mapping}
Zuk et al.\ (3.6) write the coupling in the non-irreducible form
\begin{equation}
\mu^{td}_{ij,\alpha\beta\gamma}=C_1\,\delta_{\alpha\beta}\Rhh{\gamma}+C_2\,\Rhh{\alpha}\Rhh{\beta}\Rhh{\gamma},
\label{eq:zukform}
\end{equation}
where $[\one\Rh]_{\alpha\beta\gamma}=\delta_{\alpha\beta}\Rhh{\gamma}$ in their notation.
Because $\mu^{td}$ is contracted with the symmetric traceless $E_{\beta\gamma}$, only the
part of \eqref{eq:zukform} that is symmetric and traceless in $(\beta,\gamma)$ is physical;
Zuk et al.\ deliberately record the simplest algebraic representative.  Contracting both
\eqref{eq:zukform} and the basis form \eqref{eq:tdbasis} with an arbitrary symmetric
traceless $E$ gives
\begin{equation}
[\mu^{td}\!:\!E]_\alpha
=\underbrace{C_2}_{\text{Zuk}}\,\Rhh{\alpha}(\Rh\Rh\!:\!E)+\underbrace{C_1}_{\text{Zuk}}\,\Rhh{\gamma}E_{\alpha\gamma}
=\underbrace{(x^g-2y^g)}_{\text{K\&K}}\,\Rhh{\alpha}(\Rh\Rh\!:\!E)+\underbrace{2y^g}_{\text{K\&K}}\,\Rhh{\gamma}E_{\alpha\gamma}.
\end{equation}
Matching the two independent structures $\Rhh{\alpha}(\Rh\Rh\!:\!E)$ and $\Rhh{\gamma}E_{\alpha\gamma}$
yields the linear relation
\begin{equation}
\boxed{\,x^g_{ij}=C_1+C_2,\qquad y^g_{ij}=\tfrac12\,C_1\,.}
\label{eq:mapping}
\end{equation}
This was verified numerically to machine precision against the full Cartesian tensor
\eqref{eq:zukform} in every branch, and its far-field limit reproduces the expected
$x^g\to-5a_j^3/(2R^2)$, $y^g=O(R^{-4})$.

\subsection*{Branch 1 --- non-overlapping, $R>a_i+a_j$}
Here $C_1=-\dfrac{a_j^3(5a_i^2+3a_j^2)}{3R^4}$ and
$C_2=\dfrac{5a_j^3(5a_i^2+3a_j^2-3R^2)}{6R^4}$, so \eqref{eq:mapping} gives
\begin{equation}
\boxed{\;
x^g_{ij}=\frac{a_j^3\big(5a_i^2+3a_j^2-5R^2\big)}{2R^4},
\qquad
y^g_{ij}=-\frac{a_j^3\big(5a_i^2+3a_j^2\big)}{6R^4}\;}
\qquad (R>a_i+a_j).
\end{equation}

\subsection*{Branch 2 --- overlapping, $a_>-a_<<R\le a_i+a_j$}
Zuk et al.\ write $\bm{\mu}^{td}_{ij}=\tfrac{5}{6}a_j\big[\mathscr{C}\,\one\Rh+\mathscr{D}\,\Rh\Rh\Rh\big]$,
i.e.\ $C_1=\tfrac56 a_j\mathscr{C}$, $C_2=\tfrac56 a_j\mathscr{D}$, with
\begin{align}
\mathscr{C}&=\frac{10R^6-24a_iR^5-15R^4(a_j^2-a_i^2)+(a_j-a_i)^5(a_i+5a_j)}{40\,a_i a_j R^4},\\[4pt]
\mathscr{D}&=\frac{\big((a_i-a_j)^2-R^2\big)^2\big((a_i-a_j)(a_i+5a_j)-R^2\big)}{16\,a_i a_j R^4}.
\end{align}
Then \eqref{eq:mapping} gives, compactly,
\begin{equation}
\boxed{\;
x^g_{ij}=\frac{5}{6}\,a_j\big(\mathscr{C}+\mathscr{D}\big),
\qquad
y^g_{ij}=\frac{5}{12}\,a_j\,\mathscr{C}\;}
\qquad (a_>-a_<<R\le a_i+a_j).
\end{equation}
Explicitly, in terms of $a_i,a_j,R$ (writing $P\equiv 10R^6-24a_iR^5-15R^4(a_j^2-a_i^2)+(a_j-a_i)^5(a_i+5a_j)$
and $Q\equiv\big((a_i-a_j)^2-R^2\big)^2\big((a_i-a_j)(a_i+5a_j)-R^2\big)$),
\begin{align}
y^g_{ij}&=\frac{P}{96\,a_i R^4},\\[4pt]
x^g_{ij}&=2\,y^g_{ij}+\frac{5a_j}{6}\,\mathscr{D}
=\frac{P}{48\,a_i R^4}+\frac{5\,Q}{96\,a_i R^4}
=\frac{2P+5Q}{96\,a_i R^4}.
\end{align}

\subsection*{Branch 3 --- complete overlap, $R\le a_>-a_<$}
Zuk et al.\ give $\bm{\mu}^{td}_{ij}=-\Theta(a_j-a_i)\,R\,\one\Rh$, i.e.\ $C_1=-\Theta(a_j-a_i)\,R$,
$C_2=0$, so
\begin{equation}
\boxed{\;
x^g_{ij}=-\Theta(a_j-a_i)\,R,
\qquad
y^g_{ij}=-\tfrac12\,\Theta(a_j-a_i)\,R\;}
\qquad (R\le a_>-a_<).
\end{equation}
This is non-zero only when the dipole-bearing sphere $j$ is the larger one (so that sphere
$i$ is immersed in it); otherwise the coupling vanishes.  All branches are continuous at
$R=a_i+a_j$ and $R=a_>-a_<$ (verified numerically).

\section*{Remarks for implementation}
\begin{itemize}[leftmargin=1.4em]
\item \textbf{Index/particle roles.} In $\bm{\mu}^{td}_{ij}$, particle $i$ carries the
translational velocity (Fax\'en response) and particle $j$ carries the dipole/stresslet;
the asymmetry $a_i\leftrightarrow a_j$ is genuine, so populate the transpose block
($\bm{\mu}^{dt}$, strain on $i$ from force on $j$) with the roles swapped rather than by
naive transposition.
\item \textbf{$y^g$ basis normalization.} The relations \eqref{eq:mapping} assume the
$y^g$ basis tensor $\big(\delta_{\alpha\beta}\Rhh{\gamma}+\delta_{\alpha\gamma}\Rhh{\beta}
-2\Rhh{\alpha}\Rhh{\beta}\Rhh{\gamma}\big)$ exactly as in \eqref{eq:tdbasis}.  If your code
instead defines this basis tensor with a factor $\tfrac12$ (i.e.\
$\tfrac12(\delta_{\alpha\beta}\Rhh{\gamma}+\delta_{\alpha\gamma}\Rhh{\beta})
-\Rhh{\alpha}\Rhh{\beta}\Rhh{\gamma}$), multiply the tabulated $y^g$ by $2$.  The $x^g$
basis tensor $(\Rhh{\beta}\Rhh{\gamma}-\third\delta_{\beta\gamma})\Rhh{\alpha}$ is
unambiguous.
\item \textbf{Consistency with the tt substitution rule.} For the non-overlapping
$\bm{\mu}^{tt}$ only, the different-sized result follows from the equal-sized one via
$a^2\to(a_i^2+a_j^2)/2$; this shortcut fails for all overlapping branches and for the
$rt$ and $td$ couplings, which is why the explicit forms above are needed.
\end{itemize}

\end{document}
